\chapter{Introduction}

Pneumonia is frequently evaluated in acute care using chest radiography. Chest X-rays (CXR) are typically acquired early in the clinical encounter when pulmonary infection is suspected, making them a natural anchor point for predictive modeling. Automated analysis of CXR, particularly when combined with structured clinical data, has therefore become an active area of research. However, the validity of such models depends critically on whether their inputs and evaluation protocols reflect the actual information available at the intended prediction time.

Deep learning approaches, especially convolutional neural networks, have demonstrated strong performance in image-based chest pathology classification. In parallel, models built on structured tabular data—such as emergency department (ED) triage measurements—capture physiological and demographic information available at presentation. Combining these modalities is a natural extension when both are available for the same encounter. Nonetheless, the benefit of multimodal fusion is not guaranteed: it depends on the presence of complementary signal and on the strictness of the temporal constraints imposed on the input features. In particular, distinguishing pneumonia from other radiographic abnormalities—such as pleural effusion, pulmonary edema, or non-specific opacities—is inherently challenging, as these conditions often exhibit overlapping visual patterns on chest radiographs.

In this thesis, prediction is explicitly anchored at the time of image acquisition ($t_0$). Only variables available at or before $t_0$ are permitted as inputs. This constraint addresses a key challenge in medical machine learning, namely temporal leakage, where features reflecting downstream clinical decisions or outcomes may inadvertently be used during training. While restricting inputs to a pre-imaging window reduces the amount of available information and may increase task difficulty, it ensures that the problem formulation aligns with a clearly defined and clinically meaningful decision point.

The study is based on the MIMIC-CXR-JPG (v2.1.0) dataset and structured data from MIMIC-IV and MIMIC-IV-ED. An emergency department (ED)-linked cohort of chest radiograph studies is constructed under strict linkage and filtering criteria, including the selection of a single frontal image per study with deterministic prioritization and linkage to a corresponding ED stay for each retained study. The supervised task is binary radiographic pneumonia (target T1), derived from CheXpert labels using a \texttt{u\_ignore} uncertain-label policy, where uncertain labels are excluded from training and evaluation. It is important to note that this label reflects radiographic interpretation rather than a definitive etiologic diagnosis. The final labeled dataset contains 9{,}137 studies, partitioned using a patient-level temporal split into training, validation, and test sets.

Implemented models include triage-based clinical baselines (logistic regression and gradient boosted trees), an image-only DenseNet121 model fine-tuned for pneumonia following multilabel pretraining, and an early-fusion multimodal model combining image and clinical features. All models are trained and evaluated on a shared cohort under identical temporal constraints to enable direct comparison.

Evaluation is performed on a shared held-out temporal test set of 1{,}075 studies using discrimination metrics (AUROC and AUPRC), complemented by patient-level bootstrap resampling for paired comparisons. Additional analyses include calibration assessment, decision curve analysis (DCA), and error-oriented evaluations such as CheXpert-based stratification of negative cases and gradient-based visualizations for model interpretability. DCA is interpreted in a threshold-dependent manner, reflecting variation in clinical utility across operating points.

Within this setup, image-based models achieve higher discrimination than clinical baselines (AUROC 0.746 vs 0.595), while multimodal fusion does not demonstrate a consistent improvement over the image-only model (AUROC 0.736 vs 0.746), with paired bootstrap comparisons indicating no statistically reliable advantage. Performance is further sensitive to the composition of negative cases, with reduced discrimination observed when negatives include studies with other radiographic abnormalities. This behavior highlights the inherent ambiguity of radiographic pneumonia detection and suggests that models may respond to general abnormality patterns rather than disease-specific features. These findings motivate a careful, constraint-aware interpretation of multimodal learning in this setting. This study represents a retrospective single-dataset benchmark under strict temporal constraints.

\section*{Contributions}

\begin{itemize}
    \item A reproducible pipeline for constructing an ED-linked chest radiograph cohort with prediction anchored at imaging time, including strict enforcement of pre-imaging feature availability and explicit handling of leakage-prone variables.
    
    \item A controlled comparison of clinical, image-based, and multimodal models on a shared patient-level temporal split derived from 9{,}137 labeled studies.
    
    \item A comprehensive evaluation framework incorporating discrimination metrics, patient-level bootstrap analysis, calibration assessment, and decision curve analysis to provide a multi-dimensional assessment of model performance.
    
    \item Empirical evidence that multimodal fusion does not provide a consistent performance gain over a strong image-only baseline under strict temporal constraints, highlighting the importance of problem formulation in multimodal learning.
    
    \item An analysis of error structure showing that model performance depends strongly on the definition of negative cases, with reduced discrimination in the presence of other radiographic abnormalities, reflecting the inherent ambiguity of pneumonia detection from chest X-rays.
\end{itemize}